% 03-physical.tex - Section 3: Physical Interfaces

\chapter{Physical Interfaces}
\label{sec:physical}

\section{Overview}

The NDP subsystem interfaces with the following components:
\begin{itemize}
    \item Analog Signal Processor (ASP) --- signal input
    \item Monitor and Control System (MCS) --- control interface
    \item Data recorders --- output data
    \item Network infrastructure --- data transport
\end{itemize}

\section{Signal Input from ASP}

The NDP receives conditioned analog signals from the ASP via coaxial cables. Each stand provides two polarization signals (X and Y).

\subsection{Input Specifications}

Table~\ref{tab:input-specs} lists the signal input specifications.

\begin{table}[htbp]
    \centering
    \caption{Signal Input Specifications}
    \label{tab:input-specs}
    \begin{tabular}{lll}
        \toprule
        \textbf{Parameter} & \textbf{Specification} & \textbf{Notes} \\
        \midrule
        Total inputs & 512 & 256 stands $\times$ 2 polarizations \\
        Connector Type & QMA Female & \\
        Impedance & 50~$\Omega$ & \\
        \bottomrule
    \end{tabular}
\end{table}

\section{ZCU102 F-engine Boards}

The F-engines perform analog-to-digital conversion and frequency channelization using Xilinx ZCU102 FPGA boards that each host two LWA digitizer boards.

\subsection{Board Configuration}

Table~\ref{tab:zcu-config} summarizes the ZCU102 board configuration.

\begin{table}[htbp]
    \centering
    \caption{ZCU102 Board Configuration}
    \label{tab:zcu-config}
    \begin{tabular}{ll}
        \toprule
        \textbf{Parameter} & \textbf{Value} \\
        \midrule
        Number of boards & 16 \\
        Inputs per board & 32 \\
        Stands per board & 16 \\
        Sample rate & 196~MHz \\
        ADC resolution & 8 bits \\
        \bottomrule
    \end{tabular}
\end{table}

\begin{siteNote}
Currently the LWA-NA site uses SNAP2 boards instead of the ZCU102.  The SNAP2s can host four LWA digitizer boards each.  The other parameters are the same as with the ZCU102.
\end{siteNote}

\subsection{F-engine Processing}

Each ZCU102 board performs the following processing:
\begin{enumerate}
    \item Analog-to-digital conversion at 196~MHz with 8-bit resolution
    \item Polyphase filter bank (PFB) with 4096 channels
    \item Equalization (gain adjustment per channel)
    \item Packetization and transmission to X-engines
\end{enumerate}

Table~\ref{tab:fengine-params} gives the signal processing parameters for each F-engine board.

\begin{table}[htbp]
    \centering
    \caption{F-engine Signal Processing Parameters}
    \label{tab:fengine-params}
    \begin{tabular}{ll}
        \toprule
        \textbf{Parameter} & \textbf{Value} \\
        \midrule
        FFT size & 8192 (4096 output channels) \\
        Channel bandwidth & 23.926~kHz \\
        Total bandwidth & 98~MHz \\
        Output data format & Complex 8-bit (4+4) \\
        \bottomrule
    \end{tabular}
\end{table}

\subsection{Network Interface}

Each ZCU102 board has a 40~Gbps Ethernet interface for transmitting channelized data to the X-engine servers.  Table~\ref{tab:fengine-network} lists the network configuration.

\begin{table}[htbp]
    \centering
    \caption{F-engine Network Configuration}
    \label{tab:fengine-network}
    \begin{tabular}{ll}
        \toprule
        \textbf{Parameter} & \textbf{Value} \\
        \midrule
        Interface speed & 40~Gbps \\
        Protocol & UDP \\
        Hostnames & zcu01--zcu16 \\
        \bottomrule
    \end{tabular}
\end{table}

\section{GPU Servers (X-engines)}

The X-engines perform beamforming and correlation using GPU-accelerated servers.

\subsection{Server Configuration}

Table~\ref{tab:xengine-config2} summarizes the X-engine server configuration.

\begin{table}[htbp]
    \centering
    \caption{X-engine Server Configuration}
    \label{tab:xengine-config2}
    \begin{tabular}{ll}
        \toprule
        \textbf{Parameter} & \textbf{Value} \\
        \midrule
        Number of servers & 4 \\
        Pipelines per server & 4 \\
        Hostnames & ndp1--ndp4 \\
        Processing framework & Bifrost \\
        \bottomrule
    \end{tabular}
\end{table}

\subsection{X-engine Processing}

Each X-engine pipeline performs:
\begin{itemize}
    \item Packet capture and reordering
    \item Beamforming (delay and sum with gain control)
    \item Correlation (cross-multiplication and accumulation)
    \item Triggered buffer capture
    \item Output packetization
\end{itemize}

\subsection{IPMI Interface}

Each server includes an IPMI interface for out-of-band management.  Table~\ref{tab:ipmi-config} lists the IPMI settings.

\begin{table}[htbp]
    \centering
    \caption{IPMI Configuration}
    \label{tab:ipmi-config}
    \begin{tabular}{ll}
        \toprule
        \textbf{Parameter} & \textbf{Value} \\
        \midrule
        IPMI hostnames & ndp1-ipmi -- ndp4-ipmi \\
        Capabilities & Power control, sensor monitoring \\
        \bottomrule
    \end{tabular}
\end{table}

\section{Headnode}

The headnode (hostname: \texttt{ndp}) runs the NDP control service and interfaces between the system and MCS.  Table~\ref{tab:headnode-config} lists the headnode parameters.

\begin{table}[htbp]
    \centering
    \caption{Headnode Configuration}
    \label{tab:headnode-config}
    \begin{tabular}{ll}
        \toprule
        \textbf{Parameter} & \textbf{Value} \\
        \midrule
        Hostname & ndp \\
        Control service & ndp-control \\
        MCS receive port & 1742 \\
        MCS send port & 1743 \\
        \bottomrule
    \end{tabular}
\end{table}

\section{Network Architecture}

\subsection{Data Network}

The data network carries channelized data from F-engines to X-engines:
\begin{itemize}
    \item 40~Gbps (F-engines) and 100~Gbps (X-engines) Ethernet connections
    \item Dedicated data switch
    \item Multicast or unicast UDP packets
\end{itemize}

\subsection{Control Network}

The control network provides:
\begin{itemize}
    \item SSH access to all systems
    \item IPMI management access
    \item MCS command interface
\end{itemize}

\section{Output Data Interfaces}

Output data from NDP is transmitted via UDP packets.  Table~\ref{tab:output-interfaces} summarizes the output data interfaces.

\begin{table}[htbp]
    \centering
    \caption{Output Data Interfaces}
    \label{tab:output-interfaces}
    \begin{tabular}{llll}
        \toprule
        \textbf{Output Type} & \textbf{Format} & \textbf{Protocol} & \textbf{Destination} \\
        \midrule
        DRX/PBEAM & Mark~5C & UDP & Data recorder \\
        COR & Custom & UDP & Orville Imager \\
        TBS & Custom & UDP & Data recorder \\
        TBT & Custom & UDP/File & Data recorder \textit{or} Local storage \\
        \bottomrule
    \end{tabular}
\end{table}
